\section{教学创新点}

\subsection{从"操作导向"到"思维导向"的范式转变}

传统GIS实践教学以软件操作技能为中心，学生学习大量菜单命令但缺乏对空间分析原理的深入理解\citep{汤国安2019}。本案例通过AI承担代码编写、数据处理等技术性任务，将学生从重复性操作中解放，将更多精力投入地理问题的分析与建模。

从认知负荷理论视角，GIS空间分析涉及空间信息感知、地理概念理解、分析流程记忆、软件操作执行等同时处理的认知资源。通过AI承担操作层面的认知负荷，学生认知资源可更多分配到高阶思维任务（问题分析、模式识别、决策判断），提升学习效果与思维深度。

\subsection{"双师协同"教学模式的重构}

本案例构建教师与AI协同的"双师"教学模式：

教师角色重构——从"知识传授者"到"教学设计者"：教师负责设计真实问题情境、规划AI融入节点、制定评价标准，在学生遇到困难时提供启发性提示而非直接答案，引导学生对AI输出进行批判性评估。

AI角色重构——从"辅助工具"到"技术协作者"：AI在"反事实分析"环节快速生成多种情景方案，让学生体会"同一问题、不同策略"的权衡关系；在"反思报告"环节引导学生思考技术应用的伦理边界，培养数字时代的负责任研究意识。

学生角色重构——从"被动接受"到"主动建构"："问题驱动"的AI交互方式比"指令跟随"更能促进深度学习，学生在与AI互动中学会评估输出质量，形成批判性使用AI的意识和能力。

这种人机协同模式形成"人—机—人"互动闭环：学生通过与AI互动发现问题、提出质疑；教师通过观察AI交互了解学习状态、调整教学策略；学生从教师反馈中深化对AI局限性的认识。

\subsection{人机协同素养培育：服务科研能力发展}

本案例注重培养学生的人机协同素养，为未来从事地理信息相关科研工作奠定基础。

培养AI辅助科研能力：学生在本案例中掌握的人机协同工作模式，可直接迁移到未来科研工作中。AI可辅助文献检索、数据处理、代码编写、论文写作等多个环节，显著提升科研效率。

形成批判性使用AI的意识：科研工作中使用AI工具必须保持批判性态度，本案例通过反思报告环节培养学生识别AI局限性的能力，包括数据偏见、算法黑箱、事实性错误等问题。

为GeoAI研究奠定基础：GeoAI是地理信息科学的前沿方向\citep{gai2023}，本案例让学生提前接触AI与GIS的融合应用，为其未来从事GeoAI相关研究奠定基础。